\scanpage{128}
\mspage{\handwritten{4}}

\deleted{$\operatorname{Gal}(\bar{k}/k)$, mais où \unsure{en|on} fait opérer $\operatorname{Gal}(\bar{k}/k)$ par automorphismes intérieurs. Ainsi, \deleted{\unsure{les points}} le groupe des points rationnels sur $k$ du groupe fondamental de $\operatorname{Spec}(k)$ s'identifie au centre de $\operatorname{Gal}(\bar{k}/k)$.}

Voici comment on voit l'existence du systême projectif $(\unsure{\mathbb{E}|\mathbb{B}|E}_i)$. Un couple $(\mathbb{G},z)$, avec $G\in\underline{G}$ et $z\in Z(G)$, est dit \uline{minimal} si on \inserted{ne} peut trouver un sous-groupe $G'\handwritten{\subset}\underline{G}$ de $G$, \deleted{\unsure{tel que}} $\handwritten{\ne}G$, tel que $z$ soit dans de la forme $Z(i)(z')$, avec $z'\handwritten{\in}Z(G')$ et $i:G'\handwritten{\to}G$ l'injection. \deleted{\unsure{Il}} \handwritten{O}n dit qu'un couple $(G,z)$ est \uline{majoré} par un couple $(G',z')$ si on peut trouver un homomorphisme $\inserted{u:}G'\handwritten{\to}G$ tel que $z=Z(u)(z')$. Il résulte de (iii) que tout couple $(G,z)$ est majoré par un couple minimal, et de la propriété 2) de $Z$ que si $(G',z')$ est minimal et domine $(G,z)$, alors il existe \uline{un seul} $u:G'\handwritten{\to}G$ tel que $z=Z(u)(z')$. De ceci résulte que les couples $(G,z)$ minimaux forment un systême projectif pour la relation de domination, d'ailleurs filtrant (car $(G,z)$ et $(G',z')$ sont dominés par $(G\unsure{x|\times}G',(z,z'))$ lui-même dominé par un $(G'',z'')$ minimal), c'est le systême cherché. On peut si on veut choisir un couple $(G,z)$ dans tout systême de couples \inserted{minimaux} isomorphes, de \deleted{\unsure{qui}} telle façon que l'ensemble d'indices \deleted{\unsure{donc}} $I$ devient ordonné et non seulement préordonné. (N.B. Ce genre de raisonnements formels est aussi celui qui sert en théorie des \deleted{\unsure{résidus|réseaux}} modules ...).

Je ne sais encore comment devrait être formulée la suite exacte d'homotopie, pour bien faire l'inclusion dans le groupe fondamental d'une partie infinitésimale devrait donner une théorie satisfaisante du comportement du groupe fondamental par spécialisation. J'espère que si $f:X\handwritten{\to}Y$ est un morphisme propre et séparable à fibres absolument connexes i.e. tel que $f(\underline{O}_X)=\underline{O}_Y$, $X$ étant muni \deleted{\unsure{de}} \inserted{\handwritten{pour simplifier}} d'une section \inserted{\handwritten{$s$}} sur $Y$ qui détermine des points-base sur les fibres, alors les groupes fondamentaux \uline{totaux} des fibres de $X$ forment sur $Y$ un proschéma en groupes \inserted{\handwritten{$\pi_1(X/Y,s)$}} de façon précise qu'on peut trouver un systême projectif \inserted{filtrant} $(G_i)_{i\handwritten{\in}I}$ de schémas en groupes $G_i$ 'spéciaux' sur $\unsure{\mathbb{Y}|\mathbb{S}|\mathbb{X}}$, avec des homomorphismes $G_i\handwritten{\to}G_j$ qui soient des épimorphismes de $\mathbb{Y}$-schémas (i.e. correspondant à un homomorphisme injectif de \unsure{faisceaux} cohérents sur $\mathbb{Y}$), de telle façon que les groupes fondamentaux \unsure{totaux} des fibres de $X$ se\deleted{\unsure{ront}} déduisent dudit pro-schéma en groupes par simple spécialisation. (C'est en tous cas ce que semble indiquer le cas où $X$ est un schéma abélien sur $Y$, cf plus bas).
% archive sheet number: 127
